\documentclass[11pt]{article}

\usepackage{series}
\usepackage{booktabs}

\hypersetup{
  pdftitle={Deterministic Microdynamics, Projected Laws, and the Superdeterminism Boundary},
  pdfauthor={Peter Nero},
  hidelinks
}

\newcommand{\Uu}{\mathcal U}
\newcommand{\Xx}{\mathcal X}
\newcommand{\Borel}{\mathcal B}
\newcommand{\bbP}{\mathbb P}
\newcommand{\bbE}{\mathbb E}
\newcommand{\Law}{\operatorname{Law}}

\title{Deterministic Microdynamics, Projected Laws,\\
and the Superdeterminism Boundary\\
\large Factorization, Conditional Kernels, Markov Closure, and Stabilization in MTT}

\author{Peter Nero}

\date{July 2026, Version 3}

\begin{document}
\maketitle

\begin{abstract}
A deterministic microscopic evolution need not define a closed deterministic
law for a reduced state. The exact obstruction is factorization: for a map
\(\Phi:\Uu\to\Uu\) and projection \(P:\Uu\to\Xx\), a reduced map exists
precisely when \(P\Phi\) is constant on every fiber of \(P\). Probability does
not follow from failure of this condition. After a preparation measure is
supplied, regular conditional measures on the fibers define a one-step
kernel. That kernel is Dirac exactly when the next reduced state is
conditionally almost surely constant, a weaker statement than global fiber
constancy. We give finite counterexamples proving that the same projection and
microdynamics can induce different kernels under different preparations and
that a stationary projection of deterministic dynamics need not be Markov.
Thus a stochastic or Markov effective theory requires a measure, a closure
theorem, and often a scaling limit. Existing deterministic homogenization
results provide such limits under declared assumptions; projection alone does
not. We also prove a coupled-basin contraction theorem: once a record basin is
selected, a nonnegative influence matrix with spectral radius below one gives
geometric joint stabilization. This does not select the basin. Finally,
determinism and superdeterminism are separated. In a Bell experiment,
measurement independence, Bell factorization, and operational no-signaling
are distinct conditions. The stated MTT route retains measurement
independence and upper/base locality while allowing a nonseparable state to
violate Bell factorization. Its selected state, instruments, and probability
source remain completion obligations.
\end{abstract}

\section*{Version 3 Revision Note}
\addcontentsline{toc}{section}{Version 3 Revision Note}
\begin{description}
\item[Supersedes] Version 2, released in January 2026 as
\emph{Determinism Without Superdeterminism: Projection-Induced
Stochasticity, Non-Randomness, and Cascading Stabilization in Modal Triplet
Theory}.
\item[Reason] Version 2 treated a Dirac conditional kernel as global fiber
invariance, promoted a one-step conditional law to a Markov process, and
described projection as generating stochasticity. Its Gaussian OU assumption
did not derive noise from deterministic dynamics, its selection-potential
tail claim lacked the required process-level hypotheses, and its cascading
stabilization corollary had no coupled contraction estimate. It also described
superdeterminism imprecisely.
\item[Resolution] This version separates algebraic descent, almost-sure
descent, preparation-dependent kernels, Markov closure, and homogenization. It
adds exact counterexamples, replaces the unsupported cascade by a
spectral-radius contraction theorem, and aligns the Bell discussion with the
corrected locality, measurement-independence, and factorization analysis.
\item[Retained result] The fiber-factorization criterion remains exact:
determinism upstairs need not close on a projected state. Structured
effective probability and robust downstream stabilization are possible after
their missing sources and hypotheses are supplied.
\item[Remaining boundary] MTT does not yet select a universal preparation
measure, prove Markov closure for arbitrary projected dynamics, derive all
diffusion coefficients from the physical carrier, select one outcome history,
or emit the complete Bell state-and-instrument packet.
\end{description}

\tableofcontents

\section{Four questions that must not be merged}

The word ``deterministic'' can refer to different mathematical levels. Let a
complete state \(u\in\Uu\) evolve by a map
\[
u_{n+1}=\Phi(u_n).
\]
Suppose only
\[
x_n=P(u_n)
\]
is retained. Four separate questions follow:
\begin{enumerate}
\item Does one current reduced state determine the next reduced state?
\item If not, which probability measure describes unresolved complete states?
\item Does the resulting reduced process have the Markov property?
\item In a Bell experiment, are source variables statistically independent of
later settings?
\end{enumerate}

A many-to-one \(P\) answers none of the last three questions by itself.
Noninjectivity permits hidden distinctions; it does not supply a probability
law, erase temporal memory, or correlate a source with future settings.

The corrected chain is
\[
\boxed{
\begin{gathered}
\text{deterministic complete dynamics}
\ +\ \text{projection}
\ +\ \text{preparation measure}\\
\longrightarrow
\text{conditional reduced laws}
\ +\ \text{a separate closure or limit theorem}.
\end{gathered}
}
\]

\section{Exact deterministic descent}

Let \(\Uu\) and \(\Xx\) be sets, let
\[
\Phi:\Uu\to\Uu,
\qquad
P:\Uu\to\Xx,
\]
and write \(\Xx_0=P(\Uu)\).

\begin{theorem}[Fiber-factorization criterion]
\label{thm:fiber}
There exists a map \(F:\Xx_0\to\Xx\) satisfying
\[
P\circ\Phi=F\circ P
\]
if and only if
\begin{equation}
P(u)=P(v)
\quad\Longrightarrow\quad
P(\Phi u)=P(\Phi v)
\label{eq:fiber}
\end{equation}
for every \(u,v\in\Uu\).
\end{theorem}

\begin{proof}
If the factorization exists and \(P(u)=P(v)=x\), then
\[
P(\Phi u)=F(x)=P(\Phi v).
\]
Conversely, assume equation \eqref{eq:fiber}. For \(x\in\Xx_0\), choose
any \(u\) with \(P(u)=x\) and define
\[
F(x)=P(\Phi u).
\]
The fiber condition makes this definition independent of the representative.
\end{proof}

The theorem is algebraic. If \(P\) and \(\Phi\) are measurable, measurability
of the quotient map \(F\) requires the corresponding measurable-space
hypotheses. The probabilistic version is naturally almost sure.

\begin{theorem}[Almost-sure deterministic descent]
\label{thm:as-descent}
Let \(\Uu\) and \(\Xx\) be standard Borel spaces, let \(U\) have probability
law \(\mu\), and define
\[
X=P(U),
\qquad
Y=P(\Phi U).
\]
The following are equivalent:
\begin{enumerate}
\item there is a measurable \(F:\Xx\to\Xx\) such that
\(Y=F(X)\) almost surely;
\item \(Y\) is measurable with respect to \(\sigma(X)\);
\item the regular conditional law \(\Law(Y\mid X=x)\) is a Dirac measure for
\(P_\#\mu\)-almost every \(x\).
\end{enumerate}
\end{theorem}

\begin{proof}
The equivalence of the first two statements is the Doob--Dynkin lemma for
standard Borel targets. If \(Y=F(X)\), its conditional law given \(X=x\) is
\(\delta_{F(x)}\) almost everywhere. Conversely, if the conditional law is
Dirac almost everywhere, a measurable choice of its atom gives
\(Y=F(X)\) almost surely.
\end{proof}

This corrects a central overstatement in Version 2. A Dirac conditional law
proves constancy only on the part of each fiber seen by the chosen
preparation, up to conditional null sets. It does not prove equation
\eqref{eq:fiber} at every point of every fiber.

\section{Where the reduced probability law comes from}

Let \(\mu_x\) be a regular conditional distribution of \(U\) given
\(P(U)=x\). Define
\begin{equation}
K_\mu(x,A)
=
\mu_x\!\left(
\left\{u:P(\Phi u)\in A\right\}
\right),
\qquad A\in\Borel(\Xx).
\label{eq:kernel}
\end{equation}
For \(P_\#\mu\)-almost every \(x\), this is the conditional law of the next
reduced state.

\begin{proposition}[Preparation dependence]
\label{prop:preparation}
The maps \(P\) and \(\Phi\) do not determine \(K_\mu\). Two preparation
measures can have the same reduced marginal \(P_\#\mu\) and induce different
next-step kernels.
\end{proposition}

\begin{proof}
Take
\[
\Uu=\{0,1\}\times\{0,1\},
\qquad
P(x,h)=x,
\qquad
\Phi(x,h)=(h,h).
\]
Let \(\mu\) give equal weight to \((0,0)\) and \((1,0)\), and let \(\nu\)
give equal weight to \((0,1)\) and \((1,1)\). Both projected marginals are
uniform on \(\{0,1\}\). Under \(\mu\), the next projected state is always
zero, so
\[
K_\mu(x,\cdot)=\delta_0.
\]
Under \(\nu\), it is always one, so
\[
K_\nu(x,\cdot)=\delta_1.
\]
\end{proof}

The example is small because the point is structural. Projection identifies
\((x,0)\) and \((x,1)\), but only a preparation says how the unresolved
coordinate is distributed. In MTT language, a selected upper state,
preparation protocol, invariant measure, or normal state is part of every
probability theorem.

\section{A one-step kernel is not automatically Markov}

Starting from \(U_0\sim\mu\), define
\[
U_n=\Phi^n(U_0),
\qquad
X_n=P(U_n).
\]
Equation \eqref{eq:kernel} describes a conditional one-step law for one
declared preparation. The process \((X_n)\) is Markov only if
\[
\bbP(X_{n+1}\in A\mid X_0,\ldots,X_n)
=
\bbP(X_{n+1}\in A\mid X_n)
\]
almost surely for every \(n\) and measurable \(A\). Time homogeneity requires
the right side to be the same kernel at every \(n\).

\begin{proposition}[Projected deterministic dynamics need not be Markov]
\label{prop:not-markov}
A deterministic invertible microscopic map with an invariant preparation can
have a non-Markov projected process.
\end{proposition}

\begin{proof}
Let
\[
\Uu=\{0,1,2,3\},
\qquad
\Phi(i)=i+1\pmod 4,
\]
let \(\mu\) be uniform, and set
\[
P(0)=P(1)=0,
\qquad
P(2)=P(3)=1.
\]
The stationary observed sequence is a uniformly phased version of
\[
0,0,1,1,0,0,1,1,\ldots.
\]
Conditioning only on \(X_n=0\) gives
\[
\bbP(X_{n+1}=1\mid X_n=0)=\frac12.
\]
But if \(X_{n-1}=0\) as well, the current phase is the second zero and
\[
\bbP(X_{n+1}=1\mid X_n=0,X_{n-1}=0)=1.
\]
The Markov identity fails.
\end{proof}

The hidden complete state retains phase information. Conditioning on the
current reduced value does not reconstruct the posterior distribution of
that phase. A valid Markov reduction therefore needs a sufficient-state or
lumpability theorem, a memory augmentation, or a scaling limit that removes
the memory.

\section{How deterministic dynamics can have a diffusion limit}

There is no contradiction between deterministic trajectories and a
stochastic effective limit. One begins with an ensemble of initial
conditions, a fast--slow scaling, and a functional limit theorem. Standard
deterministic homogenization makes this precise for broad classes of chaotic
systems \cite{ChevyrevEtAl2022}.

The companion MTT paper
\emph{Deterministic Projection, Diffusive Limits, and Knee-Like Threshold
Transitions} gives a concrete finite source:
\[
F_{\epsilon,s}(x,y)
=
\left(
(1-\epsilon\gamma)x
+\epsilon\alpha(s-s_0)
+\sqrt{\epsilon}\sin(2\pi y_1),
\ Ay\!\!\pmod{\mathbb Z^2}
\right),
\]
where
\[
A=
\begin{pmatrix}
2&1\\
1&1
\end{pmatrix}.
\]
For every fixed initial condition the dynamics is deterministic. With the
hidden torus coordinate sampled from its invariant measure, the projected
slow coordinate has a rigorously controlled diffusion limit, and the
Green--Kubo variance for the displayed forcing is exactly \(1/2\).

That result establishes possibility, not automatic MTT source selection. A
physical promotion still needs an intertwiner from the selected carrier,
selection of the invariant preparation, and derivation of the scale
parameters.

\subsection*{Random does not mean arbitrary}

Once a probability measure is introduced, the reduced process is random in
the mathematical sense. Its distribution may be sharply concentrated,
biased, bounded, correlated, or generated as a deterministic scaling limit.
Those properties make it structured, not ``non-random.'' Lawful stochastic
models are standard probability theory; they do not form a third logical
category between deterministic and random.

The OU and selection-potential argument from Version 2 is therefore retained
only as a conditional modeling direction. To obtain a slab-exit bound, one
must specify a valid Gaussian process, its path metric and continuity,
Lipschitz dependence of the observable on the Gaussian state, the expected
supremum or entropy bound, and the relation between threshold crossing and
physical inadmissibility. A pointwise covariance estimate does not supply all
of that.

\section{What cascading stabilization can rigorously mean}

Outcome selection and later stability are different. Suppose one
outcome-completion rule has already placed a coupled apparatus and environment
inside a declared record basin
\[
\mathcal D=D_1\times\cdots\times D_m.
\]
Let \(T=(T_1,\ldots,T_m):\mathcal D\to\mathcal D\).

\begin{theorem}[Coupled-basin contraction]
\label{thm:coupled}
Assume each \((D_i,d_i)\) is complete and there is a nonnegative matrix
\(L=(L_{ij})\) such that
\[
d_i(T_i x,T_i y)
\leq
\sum_{j=1}^m L_{ij}d_j(x_j,y_j)
\]
for all \(x,y\in\mathcal D\). If the spectral radius satisfies
\[
\rho(L)<1,
\]
then there are positive weights \(w_i\) and \(q<1\) for which \(T\) is a
contraction in
\[
d_w(x,y)
=\max_i\frac{d_i(x_i,y_i)}{w_i}.
\]
Consequently \(T\) has one fixed record in \(\mathcal D\), and every orbit in
that basin converges to it geometrically.
\end{theorem}

\begin{proof}
Choose \(q\) with \(\rho(L)<q<1\). The Neumann series
\[
w=
\sum_{k=0}^{\infty}(L/q)^k\mathbf 1
\]
converges and has strictly positive components. It obeys
\[
Lw=q(w-\mathbf 1)\leq qw.
\]
If \(d_w(x,y)=r\), then
\[
d_j(x_j,y_j)\leq r w_j,
\]
so
\[
d_i(T_i x,T_i y)
\leq r(Lw)_i
\leq qrw_i.
\]
Hence \(d_w(Tx,Ty)\leq qd_w(x,y)\). The weighted product of finitely many
complete spaces is complete, so the Banach fixed-point theorem gives the
unique fixed point and geometric convergence.
\end{proof}

This theorem supplies the missing quantitative meaning of downstream
``cascading stabilization.'' It does not say which basin is entered. Multiple
records require multiple basins, a noncontractive transition region, or an
outcome-resolved instrument. The influence matrix \(L\) and its physical
spectral-radius bound must also be derived for the apparatus under study.

\section{Determinism and superdeterminism}

Consider a Bell experiment with source variable \(\Lambda\), settings
\(A,B\), and outcomes \(R,S\).

\begin{description}
\item[Measurement independence] The source distribution is independent of
the later settings:
\[
\bbP(d\lambda\mid A=a,B=b)=\bbP(d\lambda).
\]
\item[Bell factorization] Conditional on a complete hidden variable,
\[
\bbP(r,s\mid a,b,\lambda)
=
\bbP(r\mid a,\lambda)\,
\bbP(s\mid b,\lambda).
\]
\item[Operational no-signaling] The marginal outcome law on either side does
not depend on the remote setting.
\end{description}

These conditions are not equivalent. Superdeterministic Bell models are
normally characterized by relaxing measurement independence in a
deterministic common-cause description; explicit measurement-dependent local
models exist \cite{Hall2010}. Determinism by itself neither proves nor
disproves such dependence. A joint preparation-and-setting measure must be
specified.

Bell's theorem says that measurement independence and Bell factorization
imply the Bell inequalities \cite{Bell1964}. Therefore a model reproducing
CHSH violation while retaining measurement independence must abandon Bell
factorization, even if its fundamental dynamics is deterministic. It can
still preserve algebraic microcausality and operational no-signaling.

\subsection{The corrected MTT Bell route}

The current spatial Bell paper uses the following conditional package:
\[
\boxed{
\begin{gathered}
\text{preparation-selected nonseparable state independent of later settings}\\
+\ \text{base-local descent}
+\ \text{local completely positive instruments}.
\end{gathered}
}
\]
This package can violate Bell factorization without a controllable
superluminal signal. It is not a Bell-local hidden-variable completion.
Calling it ``without superdeterminism'' is justified only when the selected
joint measure really preserves measurement independence.

The logical compatibility is established. The strict MTT source theorem is
not: selected geometry must still emit the nonseparable state, probability
functional, and detector instruments from one branch.

\section{Status ledger}

\begin{center}
\begin{tabular}{p{0.22\linewidth}p{0.67\linewidth}}
\toprule
Status & Result\\
\midrule
Exact & Global fiber criterion for deterministic descent\\
Exact & Almost-sure/Dirac conditional-law criterion\\
Exact & Same reduced marginal can yield different kernels\\
Exact & Deterministic stationary projection can be non-Markov\\
Exact & Coupled-basin stabilization when \(\rho(L)<1\)\\
Imported & Deterministic homogenization under declared fast--slow hypotheses\\
Conditional & Native MTT preparation, mixing, and diffusion coefficients\\
Conditional & Measurement-independent MTT Bell completion packet\\
Open & General Markov closure and objective one-history completion\\
\bottomrule
\end{tabular}
\end{center}

\section{Completion program}

To promote the structural framework to a physical MTT theorem:
\begin{enumerate}
\item select the complete state space and evolution on the physical branch;
\item select the retained algebra or projection;
\item derive the preparation or invariant state rather than adding it after
projection;
\item test global, almost-sure, and approximate fiber factorization;
\item prove a sufficient-state, memory-kernel, or homogenization theorem;
\item derive its error bound and coefficients from the same source;
\item derive the outcome-resolved apparatus instrument;
\item prove basin-local coupled stabilization after completion; and
\item for Bell protocols, verify measurement independence, base locality,
local instruments, and no-signaling separately.
\end{enumerate}

Each failure is informative. Exact fiber factorization gives deterministic
closure. Persistent memory demands an augmented or non-Markov theory. Failure
of a functional limit rules out the proposed diffusion. Failure of
measurement independence changes the Bell interpretation and must be
reported rather than renamed.

\section{Conclusion}

The useful insight behind ``determinism without superdeterminism'' is not that
projection magically creates randomness. It is that a deterministic complete
law need not factor through a reduced state. Once a preparation is supplied,
the unresolved fibers define conditional laws; once a closure or scaling
theorem is proved, those laws may become an effective stochastic process.

This paper now proves the exact descent and conditional-law criteria, shows
by example why preparation and memory matter, and gives a quantitative theorem
for joint record stabilization after a basin has been selected. It also fixes
the Bell language: determinism, measurement independence, Bell
factorization, and no-signaling are separate properties.

MTT has compatible exact pieces at each interface, including a deterministic
homogenization example, a restricted q79 output law, base-local Bell descent,
and basin-local contraction. What remains is their selected same-source
composition. That is the honest frontier.

% BEGIN MTT MANAGED COMPUTATIONAL EVIDENCE
\section*{Computational Evidence and Reproducibility}
The factorization and conditional-kernel results are measure-theoretic and do not depend on a Standard Model numerical fit. The open strict-upgrade ledger is a corpus boundary, not evidence for determinism, Markov closure, or freedom-of-settings assumptions.

The referenced rows are frozen to the curated results repository state identified below. Hashes are grouped in eight-character blocks for line breaking.
\begin{quote}\small
\textbf{Repository:} \url{https://github.com/PeterNero/mtt-results-repro}\\
\textbf{Commit:} \texttt{31247ebb 5c22f3fb b5443024 365433c6 ee0bff4a}\\
\textbf{Manifest:} \path{release/result_manifest.json}\\
\textbf{Manifest SHA-256:}\\
\texttt{fb399689 60b00584 631dbf53 1a708e18}\\
\texttt{ef928d6b 6d935119 c185d7f6 32b1e7cd}
\end{quote}

Tier labels are quoted verbatim from that manifest. A row used directly supports only the specific computational statement identified above; a corpus-state cross-check does not prove this paper's local theorems; and an open row is evidence of an unresolved obligation, never of closure.

\par\begingroup\dimen0=\pagegoal\advance\dimen0 by -\pagetotal\ifdim\dimen0<7\baselineskip\newpage\fi\endgroup
\paragraph{Open boundary (not evidence of closure).}
\begin{itemize}
\setlength{\itemsep}{0.25em}
\item \path{A05/strict_upgrade_ledger} (\emph{open}).\par Current 2/9 strict no-knob upgrade ledger.
\end{itemize}

No imported row changes theorem ownership or promotes a neighboring claim: all local statements retain their stated hypotheses, domains, and limitations.
% END MTT MANAGED COMPUTATIONAL EVIDENCE

\begin{thebibliography}{99}

\bibitem{Bell1964}
J.~S. Bell,
\emph{On the Einstein Podolsky Rosen Paradox},
Physics Physique Fizika \textbf{1} (1964) 195--200,
doi:10.1103/PhysicsPhysiqueFizika.1.195.

\bibitem{Hall2010}
M.~J.~W. Hall,
\emph{Local Deterministic Model of Singlet State Correlations Based on
Relaxing Measurement Independence},
Physical Review Letters \textbf{105} (2010) 250404,
doi:10.1103/PhysRevLett.105.250404.

\bibitem{ChevyrevEtAl2022}
I.~Chevyrev, P.~K. Friz, A.~Korepanov, I.~Melbourne, and H.~Zhang,
\emph{Deterministic Homogenization under Optimal Moment Assumptions for
Fast--Slow Systems. Part 2},
Annales de l'Institut Henri Poincare, Probabilites et Statistiques
\textbf{58} (2022) 1328--1350,
doi:10.1214/21-AIHP1203.

\end{thebibliography}

\end{document}
